\section{Introduction}\label{sec:introduction}

\noindent Hadronic calorimeters are typically designed around a limited set of geometric and material constraints, so the question is often not whether the detector
can be made arbitrarily larger, but whether the available depth can be used more effectively. In that setting, a natural detector-optimization question is whether
low-energy neutron detection can be improved through segment-dependent absorber and scintillator thickness variation, rather than keeping the calorimeter
uniform layer by layer.

\noindent This note studies that question for the ePIC negative eta hadronic calorimeter (nHCal). The reference geometry is a uniform 10-layer sampling calorimeter
with identical absorber and scintillator thickness in every segment. The alternative explored here is a nonuniform longitudinal parameterization in which the calorimeter
is divided into three segments, each with its own absorber and scintillator thickness. This keeps the study close to a realistic detector-design problem: the overall depth
remains in a comparable range, while the internal material distribution is allowed to vary.

\noindent Because the parameter space is continuous and multi-dimensional, a brute-force scan is inefficient. The optimization is instead approached with a
surrogate-guided workflow: simulated detector responses are used to train a regression model, and that model is then used to guide subsequent geometry proposals.

\noindent This surrogate-based approach is useful only if the geometry response contains enough structure to distinguish higher-performing configurations from statistical
fluctuations. For that reason, the study is framed not only as a detector-optimization problem, but also as a test of whether a practical machine-learning workflow can
resolve meaningful performance differences in a relatively shallow calorimeter-design landscape. The main comparison throughout the note is therefore between the uniform
nHCal baseline and the best-performing nonuniform geometries identified by the optimization procedure.
